
This chapter positions the contributions of this thesis within the broader
research literature. Because the thesis spans several distinct research areas,
the related work is organised thematically rather than chronologically. Each
section introduces the relevant body of prior work, identifies the specific
limitations or open questions that motivate the contributions of this thesis,
and explains how those contributions relate to, depart from, or extend the
existing literature. Detailed discussion of closely related work for each
individual contribution is deferred to the corresponding chapter, where it
can be presented in direct relation to the specific technical claims being made.

% ─────────────────────────────────────────────────────────────────────────────
\section{Autonomous Agents and Computer Use}
\label{sec:rw_agents}
% ─────────────────────────────────────────────────────────────────────────────

The idea of automating computer tasks through high-level instruction has a
long history in human-computer interaction research, from early work on
programming by demonstration~\cite{cypher_watch, myers_pbd} to the web
automation systems of the 2000s~\cite{leshed_CoScripter}. These early systems
operated on structured representations of the interface, such as the
application's internal state tree or the underlying HTML, and were limited to
tasks in applications they had been explicitly programmed for.

The introduction of large-scale language models shifted this paradigm
fundamentally. \cite{nakano_webgpt} demonstrated that language models could
be fine-tuned to interact with web browsers using text-based observations,
and \cite{yao_webshop} showed that language-model-based agents could complete
shopping tasks on a simulated web environment. Subsequent work extended this
to general desktop environments~\cite{OSWorld, xie_osworld},
mobile platforms~\cite{rawles_androidworld, zhang_android_eval}, and
cross-platform settings~\cite{koh_visualwebarena, zhou_webarena}.

The introduction of vision-language models (VLMs) opened a new direction:
agents that observe the interface as a screenshot rather than as structured
text. \cite{shaw_pixels} showed that a VLM-based agent could locate interface
elements from screenshots without access to the underlying accessibility tree,
and \cite{hong_cogagent} trained a specialised VLM for GUI understanding and
interaction. \cite{claude_computer_use} and \cite{openai_cua} demonstrated
capable screenshot-based agents built on large foundation models. The trend
toward vision-based observation is now dominant in the literature.

On the learning side, the dominant paradigm remains behaviour cloning from
human demonstrations~\cite{humphreys_data_driven, zheng_synapse}. Approaches
using reinforcement learning are less common due to the difficulty of
specifying automatic reward functions for open-ended computer tasks, though
recent work has explored reward models trained from human
feedback~\cite{xu_aguvis} and process reward models for step-level
supervision~\cite{pan_autonomous}.

Despite this progress, systematic evaluations consistently find that current
ACU systems fail in qualitatively similar ways across platforms and
architectures~\cite{OSWorld, koh_visualwebarena}. The specific failure
characterisation, and the argument that these failures have a structural rather
than an engineering origin, is the contribution of the survey presented in
Chapter~\ref{chap:acu}. Existing surveys of the field~\cite{durante_agent,
wang_survey_agents} provide useful taxonomies but do not develop the structural
argument that connects the observed failures to the properties of the learning
paradigm. Our survey is distinguished both by its scope, covering 87 papers
and 33 datasets, and by the analytical framing it imposes on the field.

% ─────────────────────────────────────────────────────────────────────────────
\section{Efficiency of Large Language Models}
\label{sec:rw_efficiency}
% ─────────────────────────────────────────────────────────────────────────────

The computational cost of training and serving large language models has
motivated an extensive literature on methods for reducing their resource
requirements. We distinguish three families of approach.

\textbf{Model quantisation} reduces the numerical precision of model weights
and activations, typically from 32-bit or 16-bit floating point to lower
bit-width integer or floating-point formats. Post-training quantisation methods
including GPTQ~\cite{frantar_gptq} and AWQ~\cite{lin_awq} achieve 4-bit
weight quantisation with limited average benchmark degradation.
\cite{dettmers_llmint8} introduced 8-bit quantisation with mixed-precision
decomposition, and \cite{dettmers_qlora} extended this to enable
quantised fine-tuning. Despite the widespread adoption of these methods and
their description as near-lossless in the literature, their effect on
qualitative reasoning capability, as opposed to aggregate benchmark scores,
has not been systematically studied.

\textbf{Low-rank adaptation} (LoRA)~\cite{hu_lora} fine-tunes a pre-trained
model by learning low-rank updates to its weight matrices, keeping the
original weights frozen. This reduces the number of trainable parameters by
orders of magnitude and makes fine-tuning feasible on consumer hardware.
Extensions including QLoRA~\cite{dettmers_qlora} combine quantisation with
low-rank adaptation to further reduce memory requirements.

\textbf{Knowledge distillation}~\cite{hinton_distillation} trains a smaller
student model to mimic the behaviour of a larger teacher, producing compact
models with performance closer to the teacher than a model of similar size
trained from scratch would achieve. Distillation has been applied extensively
to language models~\cite{sanh_distilbert, gu_minillm}.

The benchmark presented in Chapter~\ref{chap:contextual} evaluates seven
methods from the quantisation and low-rank adaptation families against a
common protocol that includes both quantitative benchmarks and a structured
qualitative human evaluation. The key finding, that quantisation methods
marketed as lossless produce measurable qualitative degradation in reasoning
and text quality that standard benchmarks fail to detect, is not established
by any prior work. This finding has direct implications for the ACU literature,
where reasoning quality rather than average benchmark performance is the
operationally relevant measure.

% ─────────────────────────────────────────────────────────────────────────────
\section{Compositional Generalisation}
\label{sec:rw_compositionality}
% ─────────────────────────────────────────────────────────────────────────────

The question of whether neural networks can learn compositional representations
has been debated since the earliest critiques of connectionist
systems~\cite{fodor_connectionism}. Empirical investigation began in earnest
with the SCAN benchmark~\cite{lake_scan}, which tests whether
sequence-to-sequence models can generalise a set of compositional rules to
novel combinations. Subsequent work introduced more controlled benchmarks
including COGS~\cite{kim_cogs} for language, CLEVR~\cite{johnson_clevr} for
visual reasoning, and gSCAN~\cite{ruis_benchmark} for grounded instruction
following.

A consistent finding across these benchmarks is that standard neural
architectures, including Transformers of varying sizes, fail to generalise
systematically when test combinations differ from training combinations, even
when all individual components have been seen during training. This failure
is robust: it persists across architectures, training procedures, and data
augmentation strategies~\cite{dziri_faith, ontanon_making}.

Architectural responses to the compositional generalisation gap have taken
several directions. Modular architectures~\cite{andreas_neural, akyurek_dn}
explicitly decompose the computation into reusable modules, one per semantic
primitive, which are assembled at inference time. Disentangled
representations~\cite{higgins_beta, locatello_disentangling} attempt to
separate the factors of variation in the input into independent latent
dimensions. Graph neural networks~\cite{battaglia_relational} impose a
relational inductive bias suited to structured reasoning over object-attribute
pairs.

Despite these efforts, the gap remains open. Even architectures specifically
designed with compositional inductive biases continue to fail on novel
combinations that require applying a rule to objects or settings not seen in
training~\cite{bahdanau_scan}. The COGITAO framework, presented in
Chapter~\ref{chap:contextual}, addresses a significant limitation of prior
benchmarks: the combinatorial depth of existing evaluations is shallow,
and the number of unique task rules is small relative to the scale of the
models being evaluated. COGITAO generates millions of unique task rules at
adjustable composition depth, enabling a more stringent evaluation than
existing benchmarks and revealing failure modes that shallower evaluations
do not expose.

% ─────────────────────────────────────────────────────────────────────────────
\section{Mechanistic Interpretability of Transformers}
\label{sec:rw_interpretability}
% ─────────────────────────────────────────────────────────────────────────────

Understanding what Transformer models learn and how they represent information
internally has become a significant subfield of machine learning research,
commonly described as mechanistic interpretability. The goals of this research
are both scientific, advancing understanding of how learned representations
are organised, and practical, enabling more reliable and predictable model
behaviour.

Early work focused on probing: training lightweight classifiers on top of
frozen model representations to test whether specific types of information,
such as syntactic structure or semantic roles, are linearly decodable from
the representations~\cite{ettinger_probing, tenney_bert_linguistics,
clark_bert_attention}. This work established that Transformer representations
contain substantial structured information about their input, but the probing
methodology has been criticised for conflating what information is present in
the representation with what information the model actually uses during
inference~\cite{belinkov_probing}.

Attention visualisation, which inspects the attention weight matrix $A$ as a
direct window into the model's internal computation, became popular early but
has been shown to be an unreliable guide to the model's actual information
routing~\cite{jain_attention, wiegreffe_attention}. The attention weights
reflect the query-key similarities after softmax normalisation; they do not
directly reflect the causal influence of each position on the output.

More recent mechanistic work has focused on identifying circuits, subgraphs
of the attention and MLP components that implement specific
computations~\cite{olsson_context, wang_interpretability, elhage_circuits}.
This approach has produced detailed accounts of specific behaviours, such as
induction heads that implement in-context learning and name-mover heads that
track entity references across long sequences, but does not yet provide a
general account of how information is organised across the full model.

The mathematical framework developed in Chapter~\ref{chap:selfattention}
approaches the interpretability question from a different direction. Rather
than identifying circuits or probing for specific information, we derive the
structural properties of the attention weight matrices that are predicted by
the gradient dynamics of different training objectives. This is a
theory-first approach: we prove what structure the training process must
produce, then verify it empirically. This distinguishes our contribution from
the empirical-first approaches that dominate the mechanistic interpretability
literature, and it produces theoretical guarantees rather than empirical
observations. The result, that bidirectional objectives induce symmetry and
autoregressive objectives induce directionality and column dominance, is
not established by any prior work.

Related work on the geometry of learned representations includes studies of
the rank structure of attention weight matrices~\cite{chen_bert_rank,
guo_parameter}, the emergence of linear representations of
structure~\cite{park_linear, nanda_progress}, and analyses of how different
layers contribute differently to the representation of
syntax and semantics~\cite{jawahar_what}. Our work complements these by
providing a theoretical account of why specific geometric properties arise,
grounded in the mathematics of the training process itself.

% ─────────────────────────────────────────────────────────────────────────────
\section{World Models and Predictive Representation Learning}
\label{sec:rw_worldmodels}
% ─────────────────────────────────────────────────────────────────────────────

The idea that an intelligent agent should maintain an internal model of its
environment, which it uses to simulate the consequences of its actions before
executing them, has deep roots in both cognitive science~\cite{tolman_maps}
and reinforcement learning~\cite{sutton_dyna}. The operationalisation of this
idea in modern deep learning has taken several directions.

\textbf{Model-based reinforcement learning.} Systems such as
DreamerV3~\cite{hafner_dreamer3}, MBPO~\cite{janner_when}, and
TD-MPC2~\cite{hansen_tdmpc2} learn a world model jointly with a policy and
a value function, using the world model to generate synthetic rollouts for
policy improvement. These systems achieve state-of-the-art sample efficiency
on a range of continuous control benchmarks. However, they rely on generative
world models that predict in observation space, which is computationally
expensive and allocates model capacity to predicting irrelevant surface detail.

\textbf{Latent-space world models.} An alternative approach learns the world
model in a compressed latent space rather than in observation space. RSSM, the
recurrent state space model used in DreamerV3~\cite{hafner_dreamer3}, learns
a latent transition model alongside a decoder for optional reconstruction.
\cite{hafner_mastering} demonstrated that latent-space models can achieve
superhuman performance on Atari. A key limitation of these approaches is that
the latent space is not constrained to have any particular structure, and the
transition function is a learned nonlinear function without stability guarantees.

\textbf{Self-supervised and contrastive world models.} Recent work has moved
toward learning world models through self-supervised objectives that do not
require pixel-level reconstruction. \cite{guo_byol_explore} used a
bootstrapped self-supervised objective for exploration. \cite{schwarzer_spr}
showed that training on future latent state predictions with a momentum
encoder improves sample efficiency in model-free RL. The JEPA framework of
\cite{lecun_jepa} and its visual instantiation in~\cite{assran_ijepa}
demonstrated that predicting in representation space rather than observation
space yields more semantically structured representations.

Koopman-JEPA, presented in Chapter~\ref{chap:koopman}, builds on the JEPA
framework and addresses a specific limitation that none of the above works
resolves: the exponential accumulation of prediction errors in multi-step
rollouts. Prior work on Koopman-based world models~\cite{lusch_deep_koopman,
morton_deep_variational, wehmeyer_time} has established the value of the
linear lifting approach for dynamical systems analysis, but these works use
reconstruction-based objectives, do not study the stability of the learned
transition operator systematically, and have not been applied to the
multi-step prediction problem in visual model-based RL. The combination of
the JEPA objective with progressively constrained Koopman parameterisations,
and the formal connection to discrete-time LQR via the stabilisability
analysis, constitute the novel contributions of Chapter~\ref{chap:koopman}.

% ─────────────────────────────────────────────────────────────────────────────
\section{Koopman Operator Methods in Machine Learning}
\label{sec:rw_koopman}
% ─────────────────────────────────────────────────────────────────────────────

The application of Koopman operator theory to data-driven dynamical systems
modelling began with dynamic mode decomposition (DMD)~\cite{schmid_dmd} and
its extension to nonlinear observables in extended DMD
(EDMD)~\cite{williams_edmd}. These methods compute a linear approximation of
the Koopman operator from a finite dictionary of basis functions and have been
applied to fluid dynamics~\cite{rowley_spectral}, power
systems~\cite{susuki_koopman}, and robotic locomotion~\cite{shi_deep_koopman}.

The integration of Koopman operator theory with deep learning was proposed
by~\cite{lusch_deep_koopman}, who jointly learned the lifting function and
the transition matrix using a deep autoencoder with a Koopman-structured latent
dynamics model. Subsequent work explored auxiliary networks for eigenvalue
prediction~\cite{li_learning_koopman}, variational formulations for
uncertainty quantification~\cite{morton_deep_variational}, and applications
to control~\cite{shi_deep_koopman, han_deep_koopman_control}.

Spectral constraints on the Koopman transition matrix have been studied in the
context of stability. \cite{mamakoukas_derivative} enforces stable eigenvalues
through constrained optimisation. \cite{azencot_forecasting} parameterises the
transition matrix as a product of Householder reflections to enforce
orthogonality. Our Cayley parameterisation achieves the same orthogonality
guarantee through a different algebraic construction that is both differentiable
and compatible with standard optimisation, and we provide the first analysis
connecting the spectral properties of the learned Koopman matrix to the
stabilisability precondition of discrete-time LQR.

Koopman methods have also been applied to reinforcement learning as a tool for
improving the sample efficiency of model-based planning~\cite{bevanda_koopman,
weissenbacher_koopman}. These works use the linear structure of the Koopman
latent space to enable efficient trajectory optimisation but do not address
the error accumulation problem or the stability of long-horizon rollouts.

% ─────────────────────────────────────────────────────────────────────────────
\section{Inductive Biases and Structured Learning}
\label{sec:rw_inductive}
% ─────────────────────────────────────────────────────────────────────────────

The role of inductive biases in enabling generalisation from limited data is
one of the oldest questions in machine learning, formalised by
\cite{mitchell_bias} and given a principled Bayesian treatment by
\cite{wolpert_no_free_lunch}. The resurgence of interest in inductive biases
in the deep learning era has been driven by the observation that universal
approximation ability does not imply appropriate generalisation: a model that
can fit any function is not guaranteed to fit the right one from limited data.

\textbf{Architectural inductive biases} include translational equivariance
in convolutional networks~\cite{lecun_convnet}, permutation equivariance in
graph networks~\cite{battaglia_relational}, and equivariance to broader
symmetry groups in geometric deep learning~\cite{bronstein_geometric}. These
encode prior knowledge about the invariances that are relevant for a class of
tasks, and they have been shown to dramatically improve sample efficiency
when correctly matched to the task structure.

\textbf{Domain-knowledge biases} encode task-specific prior knowledge that is
not naturally expressed as an architectural symmetry. \cite{raissi_physics}
introduced physics-informed neural networks, which add terms to the training
loss that penalise violations of known physical laws. \cite{gu_structured}
incorporated structured state-space models with specific stability properties
into sequence models, improving their ability to capture long-range
dependencies. These approaches share the principle that known structure in
the domain should be reflected in the model, rather than being left to be
discovered from data.

The Hounsfield unit masking approach presented in Chapter~\ref{chap:constraints}
belongs to the domain-knowledge bias family, encoding the biophysics of X-ray
attenuation as a hard constraint on the model's input representation. The
prior work most closely related to this contribution is the use of domain
knowledge in medical image segmentation~\cite{taghanaki_deep_semantic,
chen_deep_anatomy}, which has largely focused on anatomical structure priors
rather than physical measurement constraints.

\textbf{Retrieval-based approaches} augment a parametric model with a
non-parametric external memory, grounding the model's outputs in retrieved
documents or knowledge. Retrieval-augmented generation
(RAG)~\cite{lewis_rag} concatenates retrieved passages with the model's
input and has been shown to improve factual accuracy and reduce
hallucination. \cite{borgeaud_retro} trained a language model from scratch
with retrieval, showing that retrieval enables smaller parametric models to
match the performance of much larger ones on knowledge-intensive tasks.
The hybrid retrieval pipeline presented in Chapter~\ref{chap:constraints}
extends this line of work to the scientific literature retrieval setting, where
the informal-to-formal language gap between social media posts and academic
abstracts requires both lexical precision and dense semantic generalisation.

% ─────────────────────────────────────────────────────────────────────────────
\section{Neuro-Inspired Architectures}
\label{sec:rw_neuroinspired}
% ─────────────────────────────────────────────────────────────────────────────

The relationship between biological neural computation and artificial neural
networks has been a source of both inspiration and controversy throughout the
history of the field. The perceptron~\cite{rosenblatt_perceptron} and the
multilayer perceptron~\cite{rumelhart_learning} were explicitly motivated by
simplified models of biological neurons, but subsequent developments in deep
learning have moved away from biological plausibility as an explicit design
criterion. The success of architectures with no direct biological counterpart,
such as Transformers and convolutional networks with batch normalisation, has
led many practitioners to treat biological plausibility as an aesthetic
preference rather than a scientific constraint.

We take a different view. The specific neuro-inspired principles that motivate
the contributions of Part~\ref{part:paradigm} are not appeals to biology as
an aesthetic; they are motivated by the concrete representational properties
that biological systems possess and artificial systems currently lack:
compositional organisation, robustness to input variation, and local learning
rules that do not require a globally propagated error signal. We survey the
relevant literature with this motivation in mind.

\textbf{Capsule networks and structured assemblies.} \cite{hinton_capsules}
proposed capsule networks as a departure from the scalar activation of standard
neurons toward structured representations that encode the pose and
configuration of objects. Subsequent work including
CapsNet~\cite{sabour_capsnet} and EM routing~\cite{hinton_em_routing}
demonstrated promising results on small-scale object recognition tasks, but
scaling to natural image benchmarks has proved difficult. The Cooperative
Network Architecture presented in Chapter~\ref{chap:cna} shares the motivation
of encoding structural relationships explicitly but takes a different approach,
using dynamically assembled nets rather than capsule routing, and learning
fragments from local statistics without supervision.

\textbf{Sparse coding and dictionary learning.} \cite{olshausen_sparse}
showed that optimising for sparse representations of natural images produces
oriented Gabor-like filters closely resembling the receptive fields of
primary visual cortex neurons. This result established a computational
principle, sparsity, that explains a structural property of biological
representations from first principles. Subsequent work extended sparse coding
to hierarchical settings~\cite{lee_sparse_deep} and connected it to
deep learning through dictionary learning~\cite{sulam_multi_layer}. The net
fragments of CNA are learned from local statistical regularities in a manner
conceptually related to dictionary learning, but the assembly mechanism and
the recurrent connectivity within nets are novel.

\textbf{Biologically plausible learning rules.} A substantial literature has
investigated learning rules that are local in the sense required by biological
plausibility, including contrastive Hebbian learning~\cite{movellan_contrastive},
equilibrium propagation~\cite{scellier_equilibrium}, predictive coding with
local updates~\cite{rao_predictive, millidge_predictive_coding}, and
forward-forward algorithms~\cite{hinton_forward_forward}. These approaches
avoid the non-locality of backpropagation and, in some cases, perform
comparably on simple benchmarks. Deep Feedback Control~\cite{dfc}, discussed
in Chapter~\ref{chap:dfc}, belongs to this family, proposing a feedback
control interpretation of the layer-wise update problem.

\textbf{Liquid neural networks and continuous-time dynamics.}
\cite{hasani_liquid} introduced liquid time-constant networks, a class of
recurrent networks with dynamics governed by ordinary differential equations
inspired by the nematode nervous system. These networks exhibit compact
representations and robust extrapolation to out-of-distribution inputs,
properties that share motivation with the CNA contribution. The key difference
is that liquid networks achieve these properties through differential equation
dynamics over fixed network topologies, while CNA achieves them through the
dynamic assembly of network topology itself.

% ─────────────────────────────────────────────────────────────────────────────
\section{Positioning This Thesis}
\label{sec:rw_positioning}
% ─────────────────────────────────────────────────────────────────────────────

The preceding sections have introduced the research areas that this thesis
intersects and identified the specific open questions that motivate its
contributions. We close this chapter by making the positioning of the thesis
explicit along three dimensions.

\textbf{Diagnosis before solution.} The majority of the works surveyed above
propose solutions to specific observed failures without providing a theoretical
account of why those failures occur. This thesis takes the diagnosis seriously
as a research contribution in its own right. The mathematical framework for
self-attention geometry in Chapter~\ref{chap:selfattention} is not a stepping
stone to a proposed fix; it is an account of what the current training paradigm
produces, from which the class of appropriate responses can be derived. This
distinguishes the thesis from the bulk of the architectural and training
literature, which addresses failure modes individually and empirically.

\textbf{The constraint-paradigm distinction.} Among the neuro-inspired and
structured learning works surveyed above, there is a tendency to frame
biological plausibility and engineering tractability as competing objectives.
This thesis argues for a more nuanced view: constraint-based approaches within
the current paradigm and paradigm-level reconception are complementary, not
competing, and the appropriate choice depends on how well the structure of the
failure mode can be characterised in advance. This framing, and the taxonomy
of constraint regimes it implies, is not present in the prior literature.

\textbf{Theoretical grounding across contributions.} The contributions of this
thesis are unusually diverse in their technical domains, spanning survey
methodology, mathematical analysis of attention geometry, Koopman operator
theory applied to self-supervised learning, domain knowledge integration in
clinical segmentation, information retrieval, and neuro-inspired architecture
design. What connects them is not a shared technique but a shared argument:
that the failures of current AI systems are structurally grounded and admit
structurally grounded responses. No existing work makes this argument across
this range of domains, and the synthesis of the contributions within a single
coherent framework is itself a contribution of the thesis.